\documentclass[11pt]{article}
% Uses the canonical latex-docs semantic style in-repository. The local
% lightweight fallback is for independent compilation of this distribution.
\IfFileExists{technical-user-manual.sty}{\usepackage{technical-user-manual}}{\usepackage{../support/aaoa-standalone}}
\usepackage{mathtools}
\setDocTitle{Chapter 3: Generating Functions}
\setDocSubtitle{Guided reading, derivations and original exercises}
\setDocVersion{1.0}
\setDocStatus{Educational reference / unapproved working material}
\setDocOwner{Algorithm Engineering}
\setDocAuthor{Original companion material}
\setDocSummary{Coverage of primary book Chapter 3 with independently authored examples, proof obligations and lab transfer.}
\providecommand{\Expect}{\mathbb{E}}
\providecommand{\Prob}{\mathbb{P}}
\providecommand{\Var}{\operatorname{Var}}
\providecommand{\bigoh}{\mathcal{O}}
\begin{document}
\makedoctocpage
\section{Reading map and prerequisite check}
\noindent\textbf{Primary source sections:} 3.1 Ordinary Generating Functions; 3.2 Exponential Generating Functions; 3.3 Generating Function Solution of Recurrences; 3.4 Expanding Generating Functions; 3.5 Transformations; 3.6 Functional Equations; 3.7 Quicksort Median-of-Three Recurrence; 3.8 Counting; 3.9 Probability GFs; 3.10 Bivariate GFs; 3.11 Special Functions (printed pp. 92--140).\par
\noindent\textbf{Prerequisites:} Completed preceding mathematical-method chapters or equivalent demonstrated mastery.\par
\section{Learning outcomes}
\begin{enumerate}
\item Translate recurrences to OGF functional equations.
\item Choose OGF, EGF, PGF or bivariate series based on mathematical purpose.
\item Extract coefficients and validate low-order terms independently.
\end{enumerate}
\section{Guided technical analysis}
\subsection{OGF coefficient arithmetic}
Let $A(z)=\sum_{n\ge0}a_nz^n$ and $B(z)=\sum_{n\ge0}b_nz^n$. The product coefficient is $[z^n]AB=\sum_{k=0}^n a_kb_{n-k}$; this is a convolution, not a termwise product.
\subsection{A recurrence encoded as a rational function}
If $a_n=2a_{n-1}+1$ for $n\ge1$ and $a_0=0$, then $A(z)=2zA(z)+z/(1-z)$. Thus
\[
A(z)=\frac{z}{(1-z)(1-2z)}=\frac{1}{1-2z}-\frac{1}{1-z},\qquad a_n=2^n-1.
\]
\subsection{A probability generating function}
If $X$ is a sum of $m$ independent fair Bernoulli variables, $P(z)=\Expect[z^X]=((1+z)/2)^m$. Evaluate $P'(1)=m/2$ and $P''(1)+P'(1)-P'(1)^2=m/4$ for the mean and variance. Here independence is used for multiplication; linearity of expected values does not require it.
\section{Independent problem set}
\begin{enumerate}
\item Prove $[z^n](1-z)^{-3}=\binom{n+2}{2}$ using an algebraic identity.
\item Model the number of heads in three biased independent tosses and compute mean and variance using derivatives.
\item Give a combinatorial explanation of why EGFs include a $1/n!$ factor for labelled objects.
\end{enumerate}
\section{Laboratory transfer}
\textbf{Required laboratory:} LAB-PY02: symbolic coefficient checking and convolution tests.\par
\noindent\textbf{Minimum recorded evidence:} repository commit or source hash; test inputs/seeds; actual operation counts; predicted values; one explicit model limitation.\par
\section{Mastery gate and self-review}
Coefficients through order eight match a direct recurrence; derivation preserves constant terms and normalization.\par
\begin{enumerate}\item Can I define the random object/input exactly? \item Is my identity exact, asymptotic, or empirical? \item Which base cases and counterexamples did I check? \item What would falsify my performance hypothesis?\end{enumerate}
\section{Suggested notebook structure}
Use five named sections: model and invariants; derivation; deterministic small instances; randomized replication; discrepancy and revised hypothesis. Keep a proof separate from measurement results.

\end{document}
